\subsection*{Expected thematic clusters and coverage}

The blocks are aligned with the challenges identified in the literature
(Chapter~2) and with the system-design dimensions of RQ2 (detection, timing, addressee,
and form of intervention). The following clusters are expected in the material; they serve as
a sensitizing frame for later thematic coding, not as a closed coding scheme.

\begin{enumerate}[label=\textbf{C\arabic*}]
  \item \textbf{Goal clarity and structure.}
    Unclear or poorly communicated objectives; agenda without operational sub-goals;
    loss of a shared sense of ``on track''.
    Covered by Block~2 and Block~3.
  \item \textbf{Thematic deviation and time-to-decision.}
    Off-topic sequences, circular discussion, delayed or avoided decisions in convergent
    processes (decision-making, prioritization, problem-solving).
    Covered by Block~2 and Block~3.
  \item \textbf{Participation, hierarchy, and group dynamics.}
    Unequal speaking time, dominance, reluctance to dissent, face-saving.
    Covered by Block~2, Block~3, and Block~4.
  \item \textbf{Dynamic goals and situational uncertainty.}
    Goals that shift during the meeting; false assumptions; need to renegotiate the mandate.
    Covered by Block~2 and Block~3.
  \item \textbf{Preparation versus in-meeting facilitation.}
    What can be secured beforehand, and what can only be steered in situ.
    Covered by Block~3.
  \item \textbf{Recognizing critical moments (cues for detection).}
    Verbal, temporal, social, and remote-specific signals from which facilitators infer
    that goal orientation or decision quality is at risk. Directly informs the detection
    logic of RQ2 (\textit{What? When? Who?}).
    Covered by Block~2.
  \item \textbf{Intervention repertoire.}
    Methods, structures, and habits (e.g.\ timeboxing, roles, visualization);
    choice of addressee and channel (group vs.\ private); tone and timing.
    Covered by Block~3.
  \item \textbf{Role, agency, and limits of an AI co-facilitator.}
    Proactive versus on-demand support; essential functions; undesired automation;
    risks of unsolicited intervention (trust, control, false positives, privacy).
    Covered by Block~4.
\end{enumerate}

Removed from the earlier draft were biographical and attitudinal items (general opinion of
  online collaboration, career entry, liking the role, overall rating of online vs.\ in-person
work). They do not contribute to RQ1 or RQ2 and consume time needed for narrative accounts
of practice.